% Created 2022-09-14 mer. 17:20
% Intended LaTeX compiler: pdflatex
\documentclass{article}
\usepackage[utf8]{inputenc}
\usepackage[T1]{fontenc}
\usepackage{graphicx}
\usepackage{longtable}
\usepackage{wrapfig}
\usepackage{rotating}
\usepackage[normalem]{ulem}
\usepackage{amsmath}
\usepackage{amssymb}
\usepackage{capt-of}
\usepackage{hyperref}
\usepackage[french]{babel}
\author{Maxime Gaide}
\date{\today}
\title{Contexte et objectifs de la thèse}
\hypersetup{
 pdfauthor={Maxime Gaide},
 pdftitle={Contexte et objectifs de la thèse},
 pdfkeywords={informatique géométrique et graphique, modélisation géométrique, modélisation paramétrique, modélisation spatio-temporelle, nommage persistant, archéologie, architecture},
 pdfsubject={},
 pdfcreator={Emacs 28.1 (Org mode 9.6)}, 
 pdflang={French}}
\usepackage{biblatex}
\addbibresource{~/Documents/studies/manuscript/references/references.bib}
\begin{document}

\maketitle
Les archéologues utilisent \textbf{souvent} des logiciels de modélisation géométrique
pour élaborer des restitutions virtuelles de sites et/ou d'éléments
architecturaux.\\

Ces logiciels présentent plusieurs limitations pour le test d'hypothèse ou la
modélisation de bâtiments évoluant au fil du temps \(\to\) ils ne proposent
généralement pas des outils permettant, à la fois, de représenter les
différentes phases d'évolution au cours du temps des bâtiments, d'associer
aisément au modèles 3D des sources de données hétérogènes, de tester simplement
différentes hypothèse de restitution\\

Cela leur permet de pouvoir tester différentes hypothèses de restitution ou de
modéliser des bâtiments évoluant au fil du temps\\

Nécessité de développer un logiciel de modélisation 4D (3D+temps) permettant la
construction de bâtiments anciens pouvant évoluer dans le temps,
l'enregistrement du processus de modélisation et le rejeu automatique de ce
processus afin de pouvoir générer plusieurs variantes du même bâtiment\\

Annoter les modèles avec des données historiques hétérogènes (photographies,
notes, mesures, textures, propriétés physiques, etc.)\\

la modélisation 3D permet aux archéologues d'avoir un support visuel pour confirmer ou infirmer certaines hypothèses sur la forme ou l'époque des bâtiments.

\printbibliography
\end{document}
